\section{Discussion}
\label{sec:discussion}
Our results show that feature preservation in simplified line charts depends strongly on both the type of feature and the class of algorithm used. Local geometric features are robust, global statistical features are fragile, and algorithm performance varies widely across dataset structures. These observations highlight that no algorithm is universally best and appropriate choices depend on the intended features, data properties, and whether the goal is smoothing, downsampling, or both. Based on our analysis, we have curated the following main takeaways: 

\paragraph{\textbf{Which visual features are easy and hard to preserve?}}
Our cross-dataset grading analysis (19 algorithms $\times$ 80 datasets $\times$ similarity metrics across 12 visual features) reveals that \emph{all} features are moderately difficult to preserve, with only modest performance differences between them. Local geometric features show slight advantages. \textbf{Roughness} emerges as the easiest feature to preserve (average GPA: 1.87), with 397 of 1,520 algorithm-dataset combinations earning an A. \textbf{Slope} ($\ell_1$: 1.86 GPA, $\ell_\infty$: 1.84 GPA) and \textbf{curvature} ($\ell_1$: 1.84 GPA, $\ell_\infty$: 1.84 GPA) follow closely, with 300--338 A grades each. However, the GPA range from easiest to hardest spans only 0.08 points (4\% variation), indicating that derivative-based, local features offer only marginal advantages over global statistical features.

In contrast, global statistical features and abrupt changes are inherently more challenging to maintain. \textbf{Regimes} and \textbf{change points} (both 1.79 GPA with 339 A's), \textbf{trend} ($\ell_1$: 1.81 GPA, $\ell_\infty$: 1.79 GPA), and \textbf{level} ($\ell_1$: 1.80 GPA) are the worst-performing metrics, with only 313--339 A grades each. These quantities either integrate information across the full series (trend, level) or require detecting discrete structural breaks (regimes, change points), making them vulnerable to aggressive smoothing or downsampling that accumulates bias. Our results show that if the goal of the line chart is to convey long-term trends or identify regime shifts (e.g., climate indicators, economic trends), the data should not be simplified too aggressively, because doing so risks distorting the underlying patterns.


\paragraph{\textbf{Transformers usually win on feature preservation but do not reduce resolution.}}
Aggregating FC scores across all features and datasets, transformer algorithms such as \methodGau, \methodMean, and \methodSG consistently occupy the top global ranks. One possible reason for this is that transformers operate on the full-resolution series, modifying values while retaining all sample positions. As a result, transformers achieve 31.1\% A grades (excellent preservation) across all datasets and features---more than 10$\times$ the rate of reducers (2.8\%)---and they rarely fail catastrophically.

However, these benefits come with an important constraint: transformers do not reduce the number of points to be rendered. They lower \emph{visual complexity} but not \emph{data density}. When the goal is to reduce the number of samples, transformers cannot solve the problem. They are best viewed as high-fidelity options when preserving granularity is more important than reducing the number of points.

\paragraph{\textbf{Reducers sacrifice fidelity for dimensionality reduction, but aggregators achieve both.}}
Our family-level analysis reveals a stark performance divide. Reducers, which must discard samples to achieve compression, show the steepest cost: across all datasets and features, reducers average only 2.8\% A grades and 10.8\% excellent/good performance (GPA 1.12). Even the strongest reducer, \methodFPCS (GPA 1.80, with 242 A's across 1,680 algorithm-dataset-feature combinations), trails the top transformers (\methodGau: 3.70, \methodMean: 3.68) by nearly 2 GPA points. Other reducers perform far worse: \methodLTTB, \methodM4, and \methodMinMaxLTTB cluster around 0.92--0.96 GPA, while \methodDP achieves only 0.37 GPA despite being the most consistent algorithm overall. This reveals a fundamental limitation: geometric simplification strategies like \methodDP and \methodLTTB, which select points to minimize visual distortion, struggle to preserve statistical and frequency-domain features that depend on preserving value distributions rather than visual shape.

Surprisingly, aggregators defy this pattern. Despite also reducing point counts, aggregators achieve the highest family-level performance: 47.0\% excellent/good grades (GPA 2.28), outperforming even transformers (43.3\% excellent/good, GPA 2.13). This is because aggregators like \methodPAA and \methodASAP smooth before downsampling, stabilizing global structure and preventing the catastrophic feature loss seen in pure reducers. In effect, aggregators combine the fidelity of transformers with the compression of reducers, demonstrating that the right two-stage pipeline can preserve features while reducing resolution.

This reinforces a key design distinction in our taxonomy: transformers primarily change \emph{values} while keeping the domain fixed; reducers and aggregators primarily change the \emph{domain} (which samples survive). But the data shows that \emph{how} you change the domain matters enormously. When dimensionality reduction is required, aggregation (smooth-then-downsample) consistently outperforms direct geometric reduction, and practitioners should consider building custom aggregators by pairing feature-preserving transformers with strategic reducers.


\paragraph{\textbf{Consistency vs.\ data-dependence.}}
Grades tell us how \emph{good} an algorithm can be, but not how \emph{stable} it is across datasets. Our variance analysis (Table~\ref{tab:variance}) addresses this by measuring the fluctuation in letter grades across the 80 datasets. \methodDP is the most consistent algorithm overall (average variance 0.23), with extremely low variance on \textbf{roughness} (0.01), \textbf{slope} $\ell_1$ (0.01), and \textbf{level} $\ell_1$ (0.01). Several frequency-domain filters (\methodCheb: 0.40, \methodEllips: 0.40) and downsamplers (\methodLTTB: 0.41, \methodMinMaxLTTB: 0.43) also demonstrate low variance, making them predictable choices for automated pipelines where consistent behavior is prioritized over peak performance.

By contrast, some algorithms exhibit high variance indicating strong data dependence. Rank filters, such as \methodMin (overall variance 1.70) and \methodMax (1.25), are notable examples: they can excel on certain datasets but show dramatic performance swings across different data structures, particularly for \textbf{extrema} (Min Filter: 3.39 variance for $W_\infty$, 3.04 for $W_1$) and \textbf{spikes \& dips} (2.69 variance for $W_1$). \methodFPCS, \methodASAP, and smoothers like \methodMedian also show elevated variance (1.25--1.70 range), underscoring that adaptive strategies do not automatically translate into robust cross-dataset performance. These algorithms may achieve excellence on well-matched data but risk poor performance when assumptions are violated.

\paragraph{\textbf{Data structure can overwhelm all algorithms: the precipitation anomaly.}}
The clearest sign that data structure can dominate algorithm choice is the behavior on the six precipitation datasets \texttt{climate\_prcp}. Here, all 19 algorithms received F grades across all features: 0\% excellent/good/fair and 100\% poor.~\Cref{fig:climate_prcp} shows a line chart of \texttt{climate\_atl\_prcp} dataset. These time series are extremely sparse (70--80\% zero rainfall) with rare, high-intensity events; any window-based averaging flattens those events into zeros, downsampling risks missing them entirely, and frequency filtering treats them as high-frequency noise.

\begin{figure}[h]
    \centering
    \includegraphics[width=\linewidth]{figures/original_viz/climate_atl_prcp.pdf}
    \caption{Caption}
    \label{fig:climate_prcp}
\end{figure}


Importantly, algorithms do \emph{not} fail on temperature or wind data from the same locations and periods, confirming that the issue is not measurement quality but data structure. Event-dominated, intermittent processes simply violate the assumptions of most simplification methods. For such domains, preserving "what happened when" may fundamentally conflict with aggressive simplification.

\paragraph{\textbf{Practitioners can design custom aggregators to control the fidelity-compression trade-off.}}
Our taxonomy reveals that aggregation is simply the principled combination of two building blocks: a transformer shapes the signal to preserve desired features, then a reducer selects representative points to achieve compression. This explains why predefined aggregators like \methodASAP (GPA 2.60, ranking 5th overall) outperform all pure reducers while trailing top convolutional smoothers---they balance both objectives rather than optimizing for one.

Importantly, practitioners are not limited to predefined methods. The taxonomy enables \emph{deliberate aggregator design}: pair any feature-preserving transformer (e.g., \methodGau for roughness, \methodMedian for extrema) with any strategic reducer (e.g., \methodFPCS for perceptual fidelity, \methodLTTB for geometric simplicity) to tailor the trade-off. This compositional approach moves beyond one-size-fits-all adaptive methods toward feature-aware, application-specific pipelines.


Finally, the combination of grades and variance provides a two-dimensional view of algorithm suitability: grades capture \emph{how well} an algorithm performs for particular features and datasets, while variance captures \emph{how reliable} that performance is across domains. This perspective moves the discussion beyond "which algorithm is best?" toward "which family, or combination, is appropriate for my data, my features, and my simplification goals?" and lays the groundwork for future feature-aware, data-aware simplification methods.
